\documentclass[11pt]{article}
\usepackage[preprint]{acl}
\usepackage{times}
\usepackage{latexsym}
\usepackage{array}
\usepackage{booktabs}
\usepackage{multirow}
\usepackage{graphicx}
\usepackage{amsmath}
\usepackage{url}
\usepackage{microtype}
\usepackage[T1]{fontenc}
\usepackage[utf8]{inputenc}

\title{FIN-bench-v2 Evaluation: Comparing Frontier API Models\\
and Local Inference on Finnish Benchmarks}

\author{Janne Tompuri \\
  Department of Computer Science \\
  University of Helsinki, Finland \\}

\begin{document}
\maketitle

%% ---------------------------------------------------------------
\begin{abstract}
We evaluate twelve language model configurations on the FIN-bench-v2 Finnish
benchmark \cite{kytoniemi2025}, including four frontier API models and six local models.
Because frontier model APIs do not expose log-probabilities, we use generative
scoring throughout; our results are therefore not directly comparable to the
official FIN-bench-v2 protocol.
All models are evaluated on a fixed 1{,}146-item subset spanning 12 tasks.
Gemini 3 Flash and Gemini 3.1 Pro achieve the highest overall
normalised scores (0.919 and 0.918, respectively), though the difference
between them is small and within sampling uncertainty.
GPT-5.4 with chain-of-thought reasoning follows closely (0.907).
The best local model, Gemma 4 26B, scores 0.822, substantially narrowing
the gap to frontier models.
Chain-of-thought reasoning yields a modest aggregate gain for GPT-5.4 (+0.010)
and a substantial gain for Gemma 4 E4B (+0.059), but negligible effect for
Claude Sonnet 4.6 ($-$0.010), suggesting that the benefit is model-dependent.
Comparing Poro-8B against its base model (Llama 3.1 8B) provides a
controlled estimate of the Finnish specialisation effect: +0.137 normalised,
concentrated in language-sensitive tasks.
\end{abstract}

%% ---------------------------------------------------------------
\section{Introduction}

Finnish poses distinctive challenges for language model evaluation:
its agglutinative morphology, limited web-scale training data relative to
English, and a long-tail distribution in multilingual benchmarks
\cite{luukkonen2023fingpt,kytoniemi2025}.
The FIN-bench-v2 benchmark \cite{kytoniemi2025} provides an evaluation
suite of 12 tasks in Finnish, drawing on established benchmarks such as
ARC Challenge \cite{clark2018arc}, Belebele \cite{bandarkar2024belebele},
HellaSwag \cite{zellers2019hellaswag}, TruthfulQA \cite{lin2022truthfulqa},
SIB-200 \cite{adelani2024sib200}, and SQuAD \cite{rajpurkar2016squad},
alongside Finnish-specific tasks covering factual recall, sentiment analysis,
and ethical alignment.

To our knowledge, frontier API models have not been included in previous
FIN-bench-v2 evaluations, which focus on open-weight models with
log-probability scoring.
The effect of chain-of-thought (CoT) reasoning has also not been studied
in this Finnish benchmark context.
Additionally, Apple Silicon hardware with the MLX framework
\cite{hannun2023mlx} now enables high-quality local inference on
consumer hardware, lowering the barrier for Finnish model evaluation.

We make the following contributions:
\begin{enumerate}
  \item A direct comparison of frontier API models against local
        models on FIN-bench-v2, using a unified generative scoring protocol.
  \item A study of CoT reasoning effects across five models
        (GPT-5.4, Claude Sonnet 4.6, Gemma 4 E4B, Gemini 3 Flash, Gemini 3.1 Pro).
  \item A controlled comparison of Poro-8B against its base model
        (Llama 3.1 8B), providing an estimate of the Finnish
        specialisation effect.
  \item A reproducible evaluation pipeline for Apple Silicon, with code and
        results released publicly.
\end{enumerate}

%% ---------------------------------------------------------------
\section{Background}

\paragraph{FIN-bench-v2.}
\citet{kytoniemi2025} introduce FIN-bench-v2, a Finnish benchmark comprising
12 tasks across five categories: factual recall, commonsense reasoning,
reading comprehension, sentiment and semantic tasks, and ethical alignment.
It extends an earlier Finnish benchmark \cite{luukkonen2023fingpt} and is,
to our knowledge, the most comprehensive Finnish evaluation suite currently
available.
The official evaluation protocol uses LM Evaluation Harness with
log-probability scoring (normalised accuracy) and five prompt variants per task,
with results reported as cross-variant averages without confidence intervals.

\paragraph{Models.}
Models were selected to represent a range of scales and providers, and to
enable a direct comparison between frontier APIs and consumer-hardware local
inference.
\textit{Frontier models} (accessed via vendor API, 2026):
GPT-5.4 \cite{openai2025gpt54},
Claude Sonnet 4.6 \cite{anthropic2025claude},
Gemini 3 Flash and Gemini 3.1 Pro \cite{google2025gemini3}.
\textit{Local models} (MLX on Apple M2 Max, 32 GB unified memory):
Gemma 4 26B, Gemma 4 E4B, Gemma 3 27B
\cite{google2025gemma4};
Poro-8B (Llama-Poro-2-8B-Instruct) \cite{luukkonen2023fingpt};
Llama 3.1 8B \cite{dubey2024llama3}.
Full model identifiers and API parameters are listed in
Table~\ref{tab:model_details} (Appendix~\ref{app:models}).

\paragraph{Chain-of-thought reasoning.}
Some models support optional chain-of-thought (CoT) reasoning, while others
have it enabled by default.
We evaluate CoT effects where possible. Full API details are given in
\S\ref{sec:setup}.

%% ---------------------------------------------------------------
\section{Experimental Setup}
\label{sec:setup}

\paragraph{Dataset.}
We evaluate on a fixed subset of 12 tasks drawn from the FIN-bench-v2
dataset collection \cite{kytoniemi2025}, comprising 1{,}146 items in total.
All tasks are sourced from TurkuNLP FIN-bench-v2 datasets on Hugging Face,
with full statistics in Table~\ref{tab:dataset_stats}.
Tasks span three scoring formats.
\textit{Letter-choice MCF} (model outputs A/B/C/D): ARC Challenge FI,
Belebele FIN, GoldenSwag FI (100 items each), and TruthfulQA mc1 FI
(100 items, variable number of choices).
\textit{Word-choice MCF} (model outputs the answer word or phrase):
ScandiSent FI, SIB-200 FI (100 each), General Knowledge (70),
Analogies, Emotions, HHH Alignment (100 each), and Similarities (76).
\textit{Generative}: SQuAD FI (100 items), scored with token-level F1.
The first seven tasks (ARC through SQuAD) form the main comparable subset;
the remaining five (TruthfulQA mc1 through Similarities) extend the benchmark
to additional task types.
The subset was fixed before evaluation began and was not adjusted
based on results.

\paragraph{Scoring.}
Because frontier model APIs do not expose token log-probabilities,
we use generative scoring throughout: models produce free-form text
responses from which the predicted choice is extracted by
regex-based letter or word matching.
The extractor prioritises explicit answer lines
(``Vastaus:'', the Finnish word for ``answer'', or ``Answer:'')
before searching the full response.
This deviates from the official FIN-bench-v2 protocol
(log-probability normalised accuracy); results are therefore not directly
comparable to \citet{kytoniemi2025}.
We discuss this difference and its implications in \S\ref{sec:limitations}.

For MCF tasks we report accuracy with Wilson 95\% confidence intervals (CI).
For SQuAD FI we report token-level F1 with bootstrap CI ($B=2{,}000$).
Pairwise model comparisons use McNemar's test with Benjamini-Hochberg
false discovery rate (FDR) correction across all model pairs and tasks
\cite{mcnemar1947,benjamini1995}.
McNemar's test is appropriate because evaluation outcomes are binary per
item (correct/incorrect) and items are shared across compared models.
It tests whether two models differ in their error patterns.
We report a normalised score $\hat{s} = (s - b) / (1 - b)$, where $b$ is
the random-baseline accuracy (0.25 for four-way letter MCF,
task-specific for word-choice tasks).

\paragraph{Chain-of-thought API details.}
CoT reasoning was enabled as follows:
GPT-5.4 via the Responses API with high reasoning effort (2{,}048 output tokens);
Claude Sonnet 4.6 via extended thinking (1{,}024-token budget, 2{,}048 tokens);
Gemma 4 E4B via MLX with thinking enabled.
Gemini 3 Flash and Gemini 3.1 Pro are thinking-enabled by default
and do not support a non-thinking configuration.
Exact API identifiers and parameters are listed in
Table~\ref{tab:model_details}.

\paragraph{Temperature confound.}
Reasoning-mode APIs require temperature~1; standard runs use temperature~0.
To check whether this difference affects results, we ran two models at
temperature~1 without reasoning.
GPT-5.4 at temperature~1 yielded a normalised score of 0.892 vs.\ 0.887 at
temperature~0 ($\Delta = +0.005$); Claude Sonnet 4.6 at temperature~1 yielded
0.870 vs.\ 0.873 at temperature~0 ($\Delta = -0.003$).
Both effects are negligible, and the GPT-5.4 thinking advantage of +0.019 over
the standard run cannot be explained by temperature alone.

%% ---------------------------------------------------------------
\section{Results}
\label{sec:results}

\subsection{Overall rankings}

Table~\ref{tab:main} reports normalised mean scores across 11 tasks
(SQuAD excluded from the aggregate; see \S\ref{sec:squad});
Figure~\ref{fig:overall} shows the same results visually with 95\%~CIs.
Gemini 3 Flash and Gemini 3.1 Pro record the two highest aggregate normalised
scores (0.919 and 0.918, respectively), with confidence intervals overlapping
on most individual tasks (see Figure~\ref{fig:heatmap}).
GPT-5.4 with CoT follows at 0.907.
Among non-thinking frontier configurations,
GPT-5.4 (0.896) and Claude Sonnet 4.6 (0.880) lead,
followed by Gemma 4 26B (0.822), the top local model.

\begin{table}[t]
\centering
\small
\begin{tabular}{lrr}
\toprule
\textbf{Model} & \textbf{Norm.} & \textbf{Raw} \\
\midrule
Gemini 3 Flash               & \textbf{0.919} & 0.936 \\
Gemini 3.1 Pro               & 0.918 & 0.937 \\
GPT-5.4 (think)              & 0.907 & 0.930 \\
GPT-5.4                      & 0.896 & 0.923 \\
Claude Sonnet 4.6            & 0.880 & 0.909 \\
Claude Sonnet 4.6 (think)    & 0.869 & 0.902 \\
\midrule
Gemma 4 26B                  & 0.822 & 0.866 \\
Gemma 3 27B                  & 0.792 & 0.839 \\
Gemma 4 E4B (think)          & 0.749 & 0.809 \\
Gemma 4 E4B                  & 0.691 & 0.762 \\
Poro-8B                      & 0.633 & 0.719 \\
Llama 3.1 8B                 & 0.496 & 0.613 \\
\bottomrule
\end{tabular}
\caption{Normalised mean score (11 tasks, SQuAD excluded) and raw mean.
  Normalised: $(s - b)/(1-b)$, $b$ = random baseline.
  Top section: frontier models (API). Bottom: local models.
  Gemini models are thinking-enabled by default.}
\label{tab:main}
\end{table}

\begin{figure}[t]
  \centering
  \includegraphics[width=\columnwidth]{../figures/overall_model_comparison.pdf}
  \caption{Normalised mean score per model (11 tasks, SQuAD excluded),
    sorted by performance. Error bars show 95\%~CI based on cross-task
    score variance ($\pm 1.96 \cdot \mathrm{SE}$, $n=11$ tasks).
    The log-probability control run and the temperature~1 baseline are excluded
    (see Figure~\ref{fig:heatmap}).
    Gemini models are thinking-enabled by default (no non-thinking baseline).}
  \label{fig:overall}
\end{figure}

\subsection{Chain-of-thought reasoning effects}
\label{sec:cot}

CoT reasoning has heterogeneous effects (Table~\ref{tab:cot}).
GPT-5.4 shows a modest aggregate gain (+0.010 normalised), but with
considerable task-level variability: gains on Emotions (+0.091) and
ScandiSent (+0.080) are partly offset by losses on TruthfulQA ($-$0.026)
and Similarities ($-$0.017); no individual task difference is statistically
significant (McNemar BH-corrected).
Gemma 4 E4B shows the largest absolute gain (+0.059), with the strongest
improvements on Analogies (+0.215), ARC (+0.126), and Emotions (+0.103),
though GoldenSwag falls ($-$0.093).
Claude Sonnet 4.6 shows a small negative change ($-$0.010), within confidence
intervals on all 11 tasks.
Per-task breakdowns are in Appendix~\ref{app:per_model}.

\begin{table}[t]
\centering
\small
\begin{tabular}{lrrr}
\toprule
\textbf{Model} & \textbf{No CoT} & \textbf{CoT} & \textbf{$\Delta$} \\
\midrule
GPT-5.4           & 0.896 & 0.907 & +0.010 \\
Claude Sonnet 4.6 & 0.880 & 0.869 & $-$0.010 \\
Gemma 4 E4B       & 0.691 & 0.749 & +0.059 \\
\midrule
\multicolumn{4}{l}{\textit{Gemini: thinking-only, no non-thinking baseline}} \\
\bottomrule
\end{tabular}
\caption{Normalised mean scores with and without CoT reasoning.
  Only Gemma 4 E4B's CoT gain is statistically significant on at least
  one task (McNemar $p < 0.05$, BH-corrected: Analogies, General Knowledge).
  GPT-5.4 and Claude differences are not significant at the task level.}
\label{tab:cot}
\end{table}

\subsection{Finnish specialisation at 8B scale}
\label{sec:poro}

Poro-8B (Finnish specialist, Llama 3.1 8B base, 165B Finnish tokens)
and Llama 3.1 8B (same base, no Finnish specialisation) provide a
controlled comparison.
Poro-8B achieves a normalised score of 0.633, compared with Llama 3.1 8B at 0.496 (+0.137).
The gain is concentrated in language-sensitive tasks: GoldenSwag (+0.484),
Emotions (+0.377), ARC (+0.193), Analogies (+0.167).
World-knowledge tasks show mixed results: Poro-8B scores lower on
TruthfulQA (0.204 vs.\ 0.523), partly explained by Poro-8B's higher
abstention rate on yes/no questions in TruthfulQA.

Figure~\ref{fig:poro_vs_llama} shows the per-task differences.
In our evaluation, Finnish specialisation does not appear to close the
architectural capability gap:
Poro-8B and Gemma 4 E4B (general model, $\sim$4B active parameters)
are statistically indistinguishable on 10 out of 11 tasks after BH correction
(McNemar, $p > 0.05$).
Gemma 4 E4B's normalised average (0.691) already exceeds Poro-8B (0.633)
by 0.058, despite the E4B being a smaller active-parameter model.
This suggests that architectural improvements and large-scale pretraining
may outweigh Finnish specialisation at the 8B parameter scale,
at least in our evaluation setting.
Whether the advantage of Gemma 4 E4B stems primarily from pretraining data
volume, a tokenizer better suited to Finnish morphology, or architectural
improvements remains an open question beyond the scope of this study.

\begin{figure}[t]
  \centering
  \includegraphics[width=\columnwidth]{../figures/poro_vs_llama.pdf}
  \caption{Normalised score difference: Poro-8B minus Llama 3.1 8B
    (SQuAD excluded). Positive values indicate Finnish specialisation helps.
    Error bars show propagated 95\%~CIs.
    Gains concentrate in language-sensitive tasks (GoldenSwag, Emotions,
    ARC, Analogies). TruthfulQA shows a large reversal attributable to higher
    abstention rates in Poro-8B.}
  \label{fig:poro_vs_llama}
\end{figure}

\begin{figure*}[t]
  \centering
  \includegraphics[width=\textwidth]{../figures/task_model_heatmap.pdf}
  \caption{Per-task accuracy heatmap (12 models $\times$ 11 tasks,
    SQuAD F1 excluded). Colour scale: red = 0, yellow = 0.5, green = 1.0.
    Models ordered strongest to weakest (top to bottom).
    Per-model breakdowns with 95\%~CIs are in Appendix~\ref{app:per_model}.}
  \label{fig:heatmap}
\end{figure*}

\subsection{Per-task highlights}

\paragraph{GoldenSwag.}
Frontier models achieve near-perfect accuracy (0.954--1.000),
substantially above most local models (0.700--0.830), with the exception
of Llama 3.1 8B (0.417), which lacks Finnish specialisation.
GPT-5.4 reaches 1.000 (100/100).\footnote{Training data contamination
cannot be fully ruled out, since the original English GoldenSwag data
(WikiHow articles) may appear in GPT-5.4's pretraining corpus.
However, the Finnish version is machine-translated, and Claude Sonnet 4.6
achieves a similarly high score (0.970), which is inconsistent with
model-specific memorisation.}
\citet{kytoniemi2025} report near-random log-probability scores for their
models on GoldenSwag MCF 0-shot; our high generative scores likely reflect
strong instruction-following ability, and may not indicate superior
commonsense reasoning.

\paragraph{TruthfulQA.}
Scores range widely, from Poro-8B (0.204) to Gemini 3.1 Pro (0.957).
TruthfulQA mc1 always places the correct answer at position A by design.
Poro-8B and Llama 3.1 8B show a systematic anti-primacy bias that produces
near-random accuracy regardless of factual knowledge — a position-bias artefact
specific to generative scoring (see \S\ref{sec:limitations}).
Gemma 4 models and frontier models perform consistently well and do not exhibit
this bias.

\paragraph{HHH Alignment.}
Frontier models abstain on 18--51\% of items (safety refusals),
compared with 9--23\% for local models.
Gemini 3.1 Pro and Gemini 3 Flash show the highest abstention rates
(51\% and 33\%, respectively); GPT-5.4 and Claude Sonnet 4.6 are lower
(18--26\%).
After excluding abstentions, conditional accuracy is uniformly high
(0.766--0.970 across all models), indicating comparable alignment behaviour.
Manual inspection suggests that most abstentions are over-refusals:
models decline to answer entirely rather than selecting the lesser-harm
option, as the task format requires.
This pattern is consistent with a structural mismatch between safety-tuned
models and the HHH task design rather than an alignment failure:
models with higher abstention rates appear to apply more conservative
safety policies, declining edge-case items outright.
For practical deployments, high abstention rates indicate that the model may
decline to respond in ambiguous real-world situations, which can be a
significant constraint for Finnish-language applications requiring
broad coverage.

%% ---------------------------------------------------------------
\section{Discussion}

\paragraph{Chain-of-thought and capability headroom.}
The heterogeneous CoT effects (\S\ref{sec:cot}) are consistent with the view
that explicit reasoning is most beneficial for models with substantial headroom
on hard tasks, and yields diminishing returns for models already near the task
ceiling: Claude Sonnet 4.6 — already near ceiling on most tasks — shows
negligible change, while Gemma 4 E4B, with greater headroom, gains substantially.

\paragraph{Gemini and the thinking-only baseline.}
Gemini 3 Flash and Gemini 3.1 Pro record the highest aggregate
scores (0.919 and 0.918), but their confidence intervals overlap on most tasks,
the 0.001 aggregate difference is within sampling uncertainty, and — because
both operate exclusively in thinking mode — their scores are not directly
comparable to non-thinking frontier configurations.

\paragraph{Frontier vs.\ local gap.}
The normalised gap between the best local model (Gemma 4 26B, 0.822)
and the best frontier model (Gemini 3 Flash, 0.919) is 0.097.
Gemma 4 26B trails Claude Sonnet 4.6 (0.880) by 0.058, within
confidence intervals on some individual tasks.
On tasks where scores cluster near the ceiling (ScandiSent, HHH Alignment)
differences are negligible.
For applications where API access is restricted or latency is a concern,
Gemma 4 26B running locally may represent a viable alternative,
accepting a 0.097 normalised score gap relative to the best frontier model.

%% ---------------------------------------------------------------
\section{Limitations}
\label{sec:limitations}

\paragraph{Generative vs.\ log-probability scoring.}
We validate the scoring-method gap empirically: running Gemma 4 26B on the same 1{,}046 MCF
items with both methods yields generative mean 0.870 vs.\ log-probability
mean 0.490 (+0.380).
The gap is primarily a methodological artefact: Gemma 4's instruction-tuned
chat template adds thinking-initiation tokens (\texttt{<channel|>}) at the
last prompt position, causing the log-probability to be read from the wrong
sequence position (predicting post-thinking tokens, not letter choices).
A no-template variant — bypassing the chat template and appending an explicit
single-letter instruction — narrows the gap for letter-choice tasks
(ARC: $-57$ pp, Belebele: $-27$ pp, GoldenSwag: $-31$ pp), but does not
eliminate it.
For word-choice tasks, the gap ($-0.21$ to $-0.51$) persists because
instruction-tuned models concentrate probability mass on generated text,
not on raw token completions.
Following \citet{gu2025olmes}, we treat generative scoring as
methodologically appropriate for instruction-tuned models, and our
results measure \emph{practical task performance} rather than
token-level probability distributions.

\paragraph{Single prompt variant.}
We use prompt variant p0 only.
A sensitivity analysis on ARC Challenge FI (Gemma 4 E4B, five variants,
Figure~\ref{fig:sensitivity}) shows a spread of 0.207 in engaged
accuracy, consistent with the 0.20 spread reported for Belebele by
\citet{kytoniemi2025}.
Our Wilson CIs capture sampling uncertainty, not prompt-selection
uncertainty.
Relative rankings are likely robust (all models evaluated identically),
but absolute scores may shift with different variants.

\begin{figure}[t]
  \centering
  \includegraphics[width=\columnwidth]{../figures/arc_sensitivity.pdf}
  \caption{ARC Challenge FI accuracy across five prompt variants (p0--p4)
    for Gemma 4 E4B. Variant p0 is used throughout our evaluation.
    Spread of 0.207 illustrates prompt-selection uncertainty not captured
    by Wilson CIs. Templates follow \citet{kytoniemi2025} Appendix A.}
  \label{fig:sensitivity}
\end{figure}


\paragraph{SQuAD morphological artefact.}
\label{sec:squad}
Finnish SQuAD F1 is not comparable across model tiers.
Frontier models produce grammatically correct Finnish with appropriate
morphological case marking, which diverges from gold spans copied from
context (incorrect standalone forms).
Local models copy short spans verbatim, earning higher F1 despite
less fluent Finnish output.
For example, if the context contains \textit{Korsikalla}
(`on Corsica', Finnish adessive case), a frontier model generating
the grammatically correct standalone answer \textit{Korsika}
(`Corsica') receives near-zero F1 against the gold span, whereas a
local model copying \textit{Korsikalla} verbatim scores 1.0.
SQuAD F1 is excluded from the normalised aggregate and reported separately.
It may not serve as a reliable Finnish-competence indicator without
further analysis.
We validate the bias empirically using BERTScore \cite{zhang2020bertscore}, which measures semantic similarity
without penalising inflection mismatches.
Frontier models gain +0.28--+0.45 BERTScore F1 points relative to token-F1,
compared with +0.26--+0.34 for local models, confirming that token-F1
systematically underestimates frontier model performance on this task
(see Appendix~\ref{app:bertscore}).
Under BERTScore, Gemini 3.1 Pro leads (0.833) and the cross-tier ranking
aligns more closely with overall model quality.
LLM-as-a-judge evaluation could further reduce residual bias from
verbose frontier responses.

\paragraph{Dataset artefacts.}
Two tasks have structural properties that affect score interpretation.
TruthfulQA mc1 always places the correct answer at position A
(an intentional design choice in the original dataset \cite{lin2022truthfulqa}).
Under generative scoring this creates a strong interaction with models'
letter-choice biases: a model that systematically avoids position A will
score near chance even if its factual knowledge is adequate.
Poro-8B exemplifies this: the model produces explicit letter choices
(format ``A:~text'', ``B:~text'', etc.) but selects A in only 6\% of
responses and B in 32\%, yielding 0.204 accuracy despite producing
coherent Finnish answers.
This issue would not arise under log-probability scoring, where each option's
probability is measured independently.
SIB-200 has an unequal class distribution in our sample: the
\textit{science/technology} category accounts for 29 out of 100 items
(vs.\ a uniform expectation of 14\%), so a model that consistently
identifies category names from text would exceed the random baseline
without generalising across all categories.
Both artefacts are properties of the original datasets, not of our
evaluation pipeline.

\paragraph{Subset size.}
With $n = 70$--$100$ items per task, the smallest detectable difference
at 95\% confidence is approximately 0.07--0.09 normalised on individual tasks.
McNemar BH-correction further reduces statistical power at the multiple-comparison
level.
Some differences in Table~\ref{tab:main} are not individually significant.

%% ---------------------------------------------------------------
\section{Related Work}

\citet{kytoniemi2025} introduce FIN-bench-v2 and evaluate open-weight
Finnish and multilingual models using log-probability scoring.
\citet{gu2025olmes} analyse the effect of MCF scoring choices
(log-prob vs.\ generative, prompt format) on benchmark rankings,
motivating our methodological discussion.
\citet{luukkonen2023fingpt} describe Poro, a Finnish-specialist
language model family.

%% ---------------------------------------------------------------
\section{Conclusion}

To our knowledge, we present the first systematic evaluation of frontier
API models against local models on the FIN-bench-v2 Finnish benchmark.
Gemini 3 Flash and 3.1 Pro lead overall, closely followed by GPT-5.4 with
chain-of-thought reasoning.
The best local model (Gemma 4 26B, normalised score 0.822) substantially
narrows the gap to frontier models.
Chain-of-thought reasoning benefits GPT-5.4 and Gemma 4 E4B overall but
is negligible for Claude Sonnet 4.6.
Finnish specialisation (Poro-8B) improves over the unspecialised base model
(+0.137 normalised) on language-sensitive tasks, but does not appear to
overcome the architectural gap to newer general-purpose models at the 8B scale.
Two methodological caveats should be noted when interpreting the results.
SQuAD F1 systematically disadvantages frontier models that generate
morphologically correct Finnish rather than copying gold spans verbatim,
making cross-tier SQuAD comparisons unreliable.
TruthfulQA mc1 scores for Poro-8B and Llama 3.1 8B reflect a position-bias
artefact under generative scoring rather than actual truthfulness capability,
and should not be read as assessments of factual knowledge.
We release the evaluation pipeline, all scored outputs, and prompt templates
to enable reproducible Finnish model evaluation on Apple Silicon.

%% ---------------------------------------------------------------
\bibliography{references}

%% ---------------------------------------------------------------
\clearpage
\onecolumn
\appendix

\section{Model Details}
\label{app:models}

Table~\ref{tab:model_details} lists the exact API identifiers, backends,
and inference parameters used for each model.
Frontier models were accessed via vendor APIs in April 2026.
Local models were downloaded from Hugging Face and run via the MLX
framework \cite{hannun2023mlx} on an Apple M2 Max with 32\,GB unified memory.
Gemini models operate in thinking mode by default and do not support
a non-thinking configuration.

\begin{table*}[h]
\centering
\resizebox{\textwidth}{!}{%
\begin{tabular}{p{3.6cm}p{5.8cm}llrr}
\toprule
\textbf{Model} & \textbf{API / Model Identifier} & \textbf{Provider} & \textbf{Backend} & \textbf{Max tok.} & \textbf{Temp.} \\
\midrule
GPT-5.4              & gpt-5.4-2026-03-05                & OpenAI    & Chat API          & 512  & 0 \\
GPT-5.4 (think)      & gpt-5.4                           & OpenAI    & Responses API     & 2048 & 1 \\
Claude Sonnet 4.6    & claude-sonnet-4-6                 & Anthropic & Messages API      & 1024 & 0 \\
Claude Sonnet 4.6 (th.) & claude-sonnet-4-6            & Anthropic & Messages API      & 2048 & 1 \\
Gemini 3 Flash       & gemini-3-flash-preview            & Google    & Generative API    & 512  & 0 \\
Gemini 3.1 Pro       & gemini-3.1-pro-preview            & Google    & Generative API    & 512  & 0 \\
\midrule
Gemma 4 26B          & gemma-4-26b-it-4bit               & Local     & MLX               & 1024 & 0 \\
Gemma 4 E4B          & gemma-4-e4b-it-4bit               & Local     & MLX               & 512  & 0 \\
Gemma 4 E4B (think)  & gemma-4-e4b-it-4bit               & Local     & MLX               & 4096 & 1 \\
Gemma 3 27B          & gemma-3-27b-it-qat-4bit           & Local     & MLX               & 512  & 0 \\
Poro-8B              & Llama-Poro-2-8B-Instruct-mlx-4Bit & Local & MLX          & 512  & 0 \\
Llama 3.1 8B         & Meta-Llama-3.1-8B-Instruct-4bit   & Local & MLX          & 512  & 0 \\
\bottomrule
\end{tabular}}%
\caption{Exact model identifiers and inference parameters.
  Frontier models: top section. Local models: bottom section.
  Claude Sonnet 4.6 (think) uses extended thinking with
  budget\_tokens=1024.
  GPT-5.4 (think) uses reasoning\_effort=`high'.
  Gemma 4 E4B (think) uses enable\_thinking=True.
  All local models use 4-bit quantisation.}
\label{tab:model_details}
\end{table*}

Table~\ref{tab:dataset_stats} summarises the 12 tasks included in our evaluation:
their source datasets, sample sizes, prompt format, number of answer choices,
primary metric, and the random-chance baseline used for score normalisation.

\begin{table*}[h]
\centering
\small
\begin{tabular}{llrllrc}
\toprule
\textbf{Task} & \textbf{HuggingFace dataset (TurkuNLP/)} & \textbf{$n$} & \textbf{Format} & \textbf{Choices} & \textbf{Metric} & \textbf{Baseline} \\
\midrule
ARC-Challenge (fi)  & finbenchv2-arc-c-fi-ht                    & 100 & MCF-L & 4       & Acc.   & 0.250 \\
Belebele (fi)       & finbenchv2-belebele-fi-og                 & 100 & MCF-L & 4       & Acc.   & 0.250 \\
GoldenSwag (fi)     & finbenchv2-goldenswag-fi-ht               & 100 & MCF-L & 4       & Acc.   & 0.250 \\
ScandiSent (fi)     & finbenchv2-scandisent-fi-mini             & 100 & MCF-W & 2       & Acc.   & 0.500 \\
SIB-200 (fi)        & finbenchv2-sib-200-fi-og                  & 100 & MCF-W & 7       & Acc.   & 0.143 \\
General Knowledge   & finbenchv2-fbv1-stripped-fi-ht            &  70 & MCF-W & 4--13   & Acc.   & 0.139 \\
SQuAD (fi)          & finbenchv2-squad-strip-fi-mt              & 100 & Gen   & —       & F1     & 0.000 \\
TruthfulQA mc1 (fi) & finbenchv2-opengpt-x\_truthfulqax-fi-mt   & 100 & MCF-L & 2--11   & Acc.   & 0.219 \\
Analogies           & finbenchv2-fbv1-stripped-fi-ht            & 100 & MCF-W & 5       & Acc.   & 0.200 \\
Emotions            & finbenchv2-fbv1-stripped-fi-ht            & 100 & MCF-W & 8       & Acc.   & 0.125 \\
HHH Alignment       & finbenchv2-fbv1-stripped-fi-ht            & 100 & MCF-W & 2       & Acc.   & 0.500 \\
Similarities        & finbenchv2-fbv1-stripped-fi-ht            &  76 & MCF-W & 4--6    & Acc.   & 0.242 \\
\midrule
\textbf{Total}      &                                           & \textbf{1{,}146} & & & & \\
\bottomrule
\end{tabular}
\caption{Task summary. MCF-L = multiple-choice letter (A/B/C/D…);
  MCF-W = multiple-choice word (model outputs the answer word directly);
  Gen = generative (scored with token-level F1).
  Baseline = random-chance accuracy used for score normalisation.
  For tasks with variable choice counts (General Knowledge 4--13,
  TruthfulQA 2--11, Similarities 4--6), the baseline is the weighted mean
  of $1/n_i$ across all evaluated items.
  SIB-200 has 7 categories ($1/7 \approx 0.143$).
  The ScandiSent baseline is 0.500 (two labels: positive/negative),
  which differs from the 0.333 used in \citet{kytoniemi2025};
  the Finnish dataset split contains no neutral class.}
\label{tab:dataset_stats}
\end{table*}

%% ---------------------------------------------------------------
\clearpage
\section{Per-Model Task Breakdowns}
\label{app:per_model}

Figures~\ref{fig:app_gemini_flash}--\ref{fig:app_llama31} show
per-task accuracy (or F1 for SQuAD) with 95\% confidence intervals
for all twelve models.
Error bars show Wilson CIs for MCF tasks and bootstrap CIs ($B=2{,}000$)
for SQuAD F1.
Models are ordered as in Table~\ref{tab:main} (strongest first).

\begin{figure*}[h]
  \centering
  \begin{minipage}[t]{0.48\textwidth}
    \centering
    \includegraphics[width=\textwidth]{../figures/per_model/tasks_gemini_3_flash.pdf}
    \caption{Gemini 3 Flash, per-task scores with 95\%~CI.}
    \label{fig:app_gemini_flash}
  \end{minipage}\hfill
  \begin{minipage}[t]{0.48\textwidth}
    \centering
    \includegraphics[width=\textwidth]{../figures/per_model/tasks_gemini_3.1_pro.pdf}
    \caption{Gemini 3.1 Pro, per-task scores with 95\%~CI.}
    \label{fig:app_gemini_pro}
  \end{minipage}
\end{figure*}

\begin{figure*}[h]
  \centering
  \begin{minipage}[t]{0.48\textwidth}
    \centering
    \includegraphics[width=\textwidth]{../figures/per_model/tasks_gpt-5.4_think.pdf}
    \caption{GPT-5.4 (think), per-task scores with 95\%~CI.}
    \label{fig:app_gpt_think}
  \end{minipage}\hfill
  \begin{minipage}[t]{0.48\textwidth}
    \centering
    \includegraphics[width=\textwidth]{../figures/per_model/tasks_gpt-5.4.pdf}
    \caption{GPT-5.4, per-task scores with 95\%~CI.}
    \label{fig:app_gpt}
  \end{minipage}
\end{figure*}

\begin{figure*}[h]
  \centering
  \begin{minipage}[t]{0.48\textwidth}
    \centering
    \includegraphics[width=\textwidth]{../figures/per_model/tasks_claude_sonnet_4.6.pdf}
    \caption{Claude Sonnet 4.6, per-task scores with 95\%~CI.}
    \label{fig:app_claude}
  \end{minipage}\hfill
  \begin{minipage}[t]{0.48\textwidth}
    \centering
    \includegraphics[width=\textwidth]{../figures/per_model/tasks_claude_sonnet_4.6_think.pdf}
    \caption{Claude Sonnet 4.6 (think), per-task scores with 95\%~CI.}
    \label{fig:app_claude_think}
  \end{minipage}
\end{figure*}

\begin{figure*}[h]
  \centering
  \begin{minipage}[t]{0.48\textwidth}
    \centering
    \includegraphics[width=\textwidth]{../figures/per_model/tasks_gemma_4_26b.pdf}
    \caption{Gemma 4 26B, per-task scores with 95\%~CI.}
    \label{fig:app_gemma4_26b}
  \end{minipage}\hfill
  \begin{minipage}[t]{0.48\textwidth}
    \centering
    \includegraphics[width=\textwidth]{../figures/per_model/tasks_gemma_4_e4b_think.pdf}
    \caption{Gemma 4 E4B (think), per-task scores with 95\%~CI.}
    \label{fig:app_gemma4_e4b_think}
  \end{minipage}
\end{figure*}

\begin{figure*}[h]
  \centering
  \begin{minipage}[t]{0.48\textwidth}
    \centering
    \includegraphics[width=\textwidth]{../figures/per_model/tasks_gemma_3_27b.pdf}
    \caption{Gemma 3 27B, per-task scores with 95\%~CI.}
    \label{fig:app_gemma3_27b}
  \end{minipage}\hfill
  \begin{minipage}[t]{0.48\textwidth}
    \centering
    \includegraphics[width=\textwidth]{../figures/per_model/tasks_gemma_4_e4b.pdf}
    \caption{Gemma 4 E4B, per-task scores with 95\%~CI.}
    \label{fig:app_gemma4_e4b}
  \end{minipage}
\end{figure*}

\begin{figure*}[h]
  \centering
  \begin{minipage}[t]{0.48\textwidth}
    \centering
    \includegraphics[width=\textwidth]{../figures/per_model/tasks_poro-8b.pdf}
    \caption{Poro-8B, per-task scores with 95\%~CI.}
    \label{fig:app_poro}
  \end{minipage}\hfill
  \begin{minipage}[t]{0.48\textwidth}
    \centering
    \includegraphics[width=\textwidth]{../figures/per_model/tasks_llama_3.1_8b.pdf}
    \caption{Llama 3.1 8B, per-task scores with 95\%~CI.}
    \label{fig:app_llama31}
  \end{minipage}
\end{figure*}

%% ---------------------------------------------------------------
\clearpage
\section{SQuAD FI: Token-F1 vs.\ BERTScore}
\label{app:bertscore}

Table~\ref{tab:bertscore} compares token-level F1 and BERTScore F1
(\texttt{bert-base-multilingual-cased}) for all twelve model configurations
on SQuAD FI (100 items each).
BERTScore measures semantic similarity without penalising morphologically
correct inflection forms that differ from gold spans.
The $\Delta$ column shows the gain from switching metrics; larger values
indicate greater systematic underestimation by token-F1.

\begin{table}[h]
\centering
\small
\begin{tabular}{lrrr}
\toprule
\textbf{Model} & \textbf{Token-F1} & \textbf{BERT-F1} & \textbf{$\Delta$} \\
\midrule
Gemini 3.1 Pro               & 0.549 & \textbf{0.833} & +0.284 \\
Gemma 4 26B                  & 0.515 & 0.810 & +0.295 \\
Gemma 4 E4B (think)          & 0.510 & 0.773 & +0.263 \\
Gemma 3 27B                  & 0.463 & 0.772 & +0.309 \\
GPT-5.4 (think)              & 0.417 & 0.755 & +0.338 \\
Gemini 3 Flash               & 0.416 & 0.750 & +0.334 \\
Llama 3.1 8B                 & 0.411 & 0.751 & +0.340 \\
Poro-8B                      & 0.398 & 0.729 & +0.332 \\
GPT-5.4                      & 0.328 & 0.727 & +0.400 \\
Gemma 4 E4B                  & 0.327 & 0.715 & +0.389 \\
Claude Sonnet 4.6 (think)    & 0.252 & 0.694 & +0.442 \\
Claude Sonnet 4.6            & 0.247 & 0.694 & +0.447 \\
\bottomrule
\end{tabular}
\caption{SQuAD FI token-F1 vs.\ BERTScore F1 (\texttt{bert-base-multilingual-cased}).
  Sorted by BERTScore F1.
  Frontier models show larger $\Delta$ (+0.28--+0.45) than local models
  (+0.26--+0.34), confirming that token-F1 systematically underestimates
  frontier model performance due to morphologically correct but
  gold-span-divergent answers.
  Claude's low BERTScore despite the largest $\Delta$ is consistent with
  verbose explanatory responses that reduce both precision and recall.}
\label{tab:bertscore}
\end{table}

%% ---------------------------------------------------------------
\clearpage
\section{Prompt Templates (p0)}
\label{app:prompts}

All tasks use a single prompt formulation (p0), applied identically to all models.
Prompts are delivered as the user turn in a chat template; no system prompt
is used. Placeholders in braces are filled from the dataset fields.
We use three prompt formats depending on the task type.

\paragraph{MCF-letter (ARC-Challenge (fi), Belebele (fi), GoldenSwag (fi), TruthfulQA mc1 (fi)).}
The model is asked to output a single letter (A, B, C, …).
The example below shows the ARC-Challenge p0 template:

\vspace{4pt}
{\setlength{\fboxrule}{0.4pt}\setlength{\fboxsep}{6pt}%
\noindent\fbox{\begin{minipage}{\dimexpr\linewidth-2\fboxsep-2\fboxrule\relax}
\small\ttfamily
Mikä on paras vastaus kysymykseen \{question\}?\\
\phantom{x}A \{choice\_A\}\\
\phantom{x}B \{choice\_B\}\\
\phantom{x}C \{choice\_C\}\\
\phantom{x}D \{choice\_D\}\\
Vastaus:
\end{minipage}}}
\vspace{4pt}

GoldenSwag prepends the story context before the choices.
Belebele prepends the passage and question above the choices.
TruthfulQA mc1 prepends the instruction
``\textit{Vastaa seuraavaan kysymykseen. Oikeita vastauksia on vain yksi.}''
and labels a variable number of answer options (typically four to eight, A onwards).

\paragraph{MCF-word (ScandiSent (fi), SIB-200 (fi), General Knowledge,
  Analogies, Emotions, HHH Alignment, Similarities).}
The model is asked to output one of the listed answer words or phrases.
The example below shows the ScandiSent p0 template:

\vspace{4pt}
{\setlength{\fboxrule}{0.4pt}\setlength{\fboxsep}{6pt}%
\noindent\fbox{\begin{minipage}{\dimexpr\linewidth-2\fboxsep-2\fboxrule\relax}
\small\ttfamily
Onko tekstissä esiintyvä tunne ``positiivinen'' vai ``negatiivinen''?\\
Teksti: \{text\}\\
Tunne:
\end{minipage}}}
\vspace{4pt}

SIB-200 lists the seven topic categories inline in the prompt.
General Knowledge, Analogies, HHH Alignment, and Similarities list
choices as ``\texttt{vaihtoehto: \{choice\}}'' on separate lines.
Emotions lists eight emotion labels inline:
``\textit{hämmästys, ilo, inho, luottamus, odotus, pelko, suru, suuttumus}''.

\paragraph{Generative (SQuAD fi).}
The model generates a free-form answer, scored with token-level F1
against all reference answers:

\vspace{4pt}
{\setlength{\fboxrule}{0.4pt}\setlength{\fboxsep}{6pt}%
\noindent\fbox{\begin{minipage}{\dimexpr\linewidth-2\fboxsep-2\fboxrule\relax}
\small\ttfamily
Otsikko: \{title\}\\[2pt]
Teksti: \{context\}\\[2pt]
Kysymys: \{question\}\\
Vastaus:
\end{minipage}}}
\vspace{4pt}

\end{document}
